\part*{Part III --- MMVSD}
\addcontentsline{toc}{part}{Part III --- MMVSD}

\section{The MMVSD model}\label{sec:mmvsdmodel}

\subsection{Why one envelope is not enough}

In kFDS the same envelope multiplies all frequencies. When the windows are long compared with the
response times, every oscillator follows the envelope and, by~\eqref{eq:kmod}, every response has the
same kurtosis $\kappa_z = 3\tavg{e^4}/\tavg{e^2}^2$, whatever $\fn$. Measured vibration rarely
behaves like this. A vehicle combines a road input whose level changes every few seconds (leptokurtic
responses at low $\fn$), an engine that excites a middle band stationarily (Gaussian responses there)
and impacts that excite high frequencies in short bursts (strongly leptokurtic responses at high
$\fn$). The kurtosis of the responses is a function $\kappa_z(\fn)$, and a test that reproduces the FDS
with a different function reproduces the damage with a different mix of cycles
(Section~\ref{sec:kzdamage}).

MMVSD (Moment-Matched Variable Spectral Density) prescribes both functions: the FDS $D_t(\fn)$ and the
kurtosis of the responses $\kappa_t(\fn)$, measured on the reference (or combined over a mission
with~\eqref{eq:kzmix}). It extends the idea of a variable spectral density, a PSD whose shape changes
with time, to a continuous random field designed so that the two families of moments of the
responses, their damage and their kurtosis, match the target.

\subsection{An evolutionary spectrum}\label{sec:evol}

The test signal is described by a \emph{local spectrum} that changes with time:
\begin{result}[MMVSD signal model]
\begin{equation}\label{eq:evol}
  G(f,t) = P(f)\,m(f,t),\qquad m(f,t)\ge0,\qquad \tavg{m(f,\cdot)} = 1 \ \text{for every } f,
\end{equation}
and, given the field $m$, the signal is a Gaussian process with this local spectrum.
\end{result}
The time average of the local spectrum is $\tavg{G(f,\cdot)} = P(f)\tavg{m(f,\cdot)} = P(f)$: the PSD
of the test is exactly $P$, whatever the modulation. The \emph{modulation field} $m$ says when and in
which bands the energy is concentrated; its distribution sets how leptokurtic the responses are.

$P$ is parametrised as a smooth correction of a \emph{shape prior} $\Pi(f)$, the PSD of the reference
at its full resolution:
\begin{equation}\label{eq:prior}
  P(f) = \Pi(f)\,p(f),\qquad p(f) = \sum_j p_j\,\phi_j(f),
\end{equation}
where $\phi_j$ are the hat functions of the nodes (Section~\ref{sec:lalanne}) extended flat over $3/Q$
beyond the first and last node. The peaks, notches and band edges of the reference are thus kept at
full resolution, and the unknowns are the $n$ nodal factors $p_j$. Outside the band of the grid the
test has no energy. The modulation is likewise described by its nodal values $m_j(t)$.

\subsection{Local moments of the responses}\label{sec:operators}

For a stationary input the spectral moments of the response of oscillator $i$
are~\eqref{eq:moments}. For the local spectrum~\eqref{eq:evol}, with $P$ and $m$ written through their
nodal values, the moments that the response would have if the local spectrum at time $t$ lasted for
ever are
\begin{equation}\label{eq:Sinf}
  S^{\infty}_{k,i}(t) = \sum_j W^{(k)}_{ij}\,p_j\,m_j(t),
  \qquad
  W^{(k)}_{ij} = \int f^k\,|H_{d,i}(f)|^2\,|H_{\mathrm{pre}}(f)|^4\,\Pi(f)\,\phi_j(f)\,\dd f,
\end{equation}
for $k = 0,1,2,4$. The \emph{moment operators} $W^{(k)}$ generalise Lalanne's operator ($W^{(0)}$ with
$\Pi=1$, the analogue response and no edge extension is the matrix $W$ of~\eqref{eq:W}). MMVSD uses
the response $H_d$ of the digital oscillator~\eqref{eq:Hd} and the pre-filter of
Section~\ref{sec:digital}, so that the model describes exactly the quantities that the measurement
computes. The integrals are evaluated by the trapezoidal rule with 128 intervals per segment between
nodes. The model assumes that within the band of an oscillator the local spectrum is the
interpolation of its nodal values, i.e.\ that the modulation varies slowly along the frequency axis
compared with the bandwidth $\fn/Q$.

The response does not follow the local spectrum instantly but with the envelope law~\eqref{eq:tau}.
Applying the low-pass filter $F_i$ of time constant $\tau_i = Q/(2\pi f_{\mathrm{n},i})$ to every
moment gives the \emph{local moments}
\begin{result}[Local moments of the responses]
\begin{equation}\label{eq:S}
  S_{k,i}(t) = F_i\!\left[\,\sum_j W^{(k)}_{ij}\,p_j\,m_j(\cdot)\right]\!(t),\qquad
  F_i[y](t) = \int_0^\infty \frac{1}{\tau_i}e^{-u/\tau_i}\,y(t-u)\,\dd u .
\end{equation}
\end{result}
$s_i(t) = S_{0,i}(t)$ is the local variance of the response of oscillator $i$. On the discrete time
grid of the model (step $\delta$, see Section~\ref{sec:frames}) the filter becomes the recursion
$y^F_l = a_iy^F_{l-1} + (1-a_i)y_l$ with $a_i = e^{-\delta/\tau_i}$, whose gain at zero frequency is
$(1-a_i)/(1-a_i) = 1$. The test profile is periodic (it may be played several times), so the
recursion is applied circularly, by multiplying the discrete Fourier transform of $y$ by
$(1-a_i)/(1-a_ie^{-2\pi\ii q/N_t})$ at the frequency index $q$ of the $N_t$ points of the grid.

\subsection{Kurtosis of the responses}\label{sec:kurtmodel}

Given $m$, the response of oscillator $i$ is Gaussian with the local variance $s_i(t)$. By the last
form of~\eqref{eq:kmod},
\begin{result}[Response kurtosis in the model]
\begin{equation}\label{eq:kappamodel}
  \kappa_i = 3\,\frac{\tavg{s_i^2}}{\tavg{s_i}^2}.
\end{equation}
\end{result}
Note that $\tavg{s_i} = \sum_jW^{(0)}_{ij}p_j\tavg{m_j} = \sum_jW^{(0)}_{ij}p_j$ does not depend on
the modulation (the filter keeps the mean and $\tavg{m_j}=1$): the modulation changes the kurtosis
and the damage, never the mean response variance.

\subsection{Damage of the responses}\label{sec:damagemodel}

Over a stretch of time short compared with the variations of $s_i$ but containing many cycles, the
response is a stationary Gaussian process with the moments $S_{k,i}(t)$. Its damage rate is the
expected damage per second of such a process: Dirlik's rate (Appendix~\ref{app:dirlik}), a function
$d_{\mathrm{Dir}}(S_0,S_1,S_2,S_4)$ that reduces to the narrow-band rate $\nu_0^+(2S_0)^{b/2}
\Gamma(1+b/2)$ for narrow bands, multiplied by the sampling factor~\eqref{eq:sampling} evaluated at
the mean crossing rate $\nu_0^+ = \sqrt{S_2/S_0}$ (the cycles that dominate the damage are those of
the zero crossings):
\begin{result}[Damage in the model]
\begin{equation}\label{eq:Dmodel}
  D_i = \frac{K^b}{C}\,T\,\Bigl\langle d\bigl(S_{0,i},S_{1,i},S_{2,i},S_{4,i}\bigr)\Bigr\rangle ,
  \qquad
  d = d_{\mathrm{Dir}}(S_0,S_1,S_2,S_4)\;c_b\!\left(\frac{\sqrt{S_2/S_0}}{f_s}\right).
\end{equation}
\end{result}
For a stationary Gaussian signal ($m\equiv1$) and a narrow band this is Lalanne's
formula~\eqref{eq:dgauss}, corrected for the digital oscillator and the sampling. The program also
offers the plain narrow-band (Rice) rate instead of Dirlik's.

How well does the model describe a realisation? For a stationary Gaussian profile synthesised by the
program at $f_s = 500$\,Hz with nodes up to 150\,Hz ($Q=10$, $b=8$), the measured response variances
agree with $S_0$ within 2--7\,\% at all nodes, and the ratio of the rainflow damage to~\eqref{eq:Dmodel}
averages about 0.9 (the remaining difference is the scatter of the largest cycles of a 60\,s
realisation, whose median is below its mean). With the analogue operators and without the sampling
factor, the same ratio drops to 0.15 at 113\,Hz.

\subsection{The modulation field}\label{sec:field}

The field must have, at every node, a prescribed depth of modulation, and between nodes and over time
the correlations that describe the bursts of the reference: how long a burst lasts (\emph{burst
duration} $T_c$, possibly different at every node) and how wide a band it covers (\emph{burst width}
$B$, in octaves). It is built in three steps on a grid of $K$ time frames (spacing $h_t$, the
\emph{hop}).

\emph{1. A latent Gaussian field.} Draw independent standard normal numbers $\xi_{j\ell}$ for every
node $j$ and frame $\ell$. Smooth every row in time with a Gaussian kernel of standard deviation
$T_{c,j}/h_t$ frames (circularly, by multiplying its discrete Fourier transform by
$\exp(-2\pi^2\varsigma_j^2\nu^2)$, $\nu$ in cycles per frame, $\varsigma_j = T_{c,j}/h_t$), standardise
each row to unit variance, then mix the rows with the Gaussian kernel
$\exp\bigl(-\tfrac12(\log_2f_{\mathrm{n},j}-\log_2f_{\mathrm{n},j'})^2/B^2\bigr)$ along the logarithmic
frequency axis, and standardise again. The result $Z_{j\ell}$ is a Gaussian field with zero mean, unit
variance at every node, and the prescribed correlations. A Gaussian kernel of standard deviation
$\varsigma$ applied to white noise gives the correlation $e^{-\lambda^2/(4\varsigma^2)}$ at the lag
$\lambda$ (the autocorrelation of the kernel is a Gaussian of variance $2\varsigma^2$); the correlation
falls to $1/e$ at $\lambda = 2\varsigma$. The settings $T_c$ and $B$ are therefore half the $1/e$
widths of the correlations.

\emph{2. Ranks.} In every row, replace the values by their ranks $r_{j\ell}\in\{0,\dots,K-1\}$ over
time. The ranks keep the information of when the field is high or low, which is all that the copula
(the dependence structure) needs.

\emph{3. Marginal.} The Gamma distribution of shape $a>0$ and scale $1/a$ has the density
$p(m) = a^a m^{a-1}e^{-am}/\Gamma(a)$ for $m>0$. With the substitution $u = am$, its moments are
$\E[m^k] = \frac{1}{a^k\Gamma(a)}\int_0^\infty u^{a+k-1}e^{-u}\dd u = \frac{\Gamma(a+k)}{a^k\Gamma(a)}$;
since $\Gamma(s+1) = s\Gamma(s)$ (integrate by parts), $\E[m] = \frac{a\Gamma(a)}{a\Gamma(a)} = 1$ and
$\E[m^2] = \frac{(a+1)a}{a^2} = 1+\frac1a$, so the variance is $1/a$. Map the ranks to the quantiles of
this distribution with $a_j = 3/e_j$ (mean 1, variance $e_j/3$):
\begin{equation}\label{eq:gammafield}
  m_{j\ell} = F^{-1}_{a_j}\!\left(\frac{r_{j\ell}+\tfrac12}{K}\right),
\end{equation}
and divide every row by its mean so that $\tavg{m_j} = 1$ exactly. Because every row uses the same set
of $K$ quantiles, the sample moments of $m_j$ depend only on $e_j$: there is no sampling noise in
$\tavg{m_j^2}$ or $\tavg{m_j^{b/2}}$ (a \emph{stratified} copula). The quantiles of the Gamma
distribution are tabulated for 200 shapes and interpolated, which makes the field cheap to rebuild
during the design.

\emph{Meaning of the depth $e_j$.} A Gaussian signal modulated directly by $m$ (local variance
proportional to $m$) has, by~\eqref{eq:kmod}, the kurtosis $3\tavg{m^2}/\tavg{m}^2 =
3(1+\Var m) = 3 + e$; $e$ is the \emph{excess kurtosis of the modulation}. The kurtosis of a
response is lower when the bursts are short compared with $\tau_i$, by~\eqref{eq:papoulis}; the design
(Section~\ref{sec:design}) finds the $e_j$ that give the target $\kappa_t$ through the full
model~\eqref{eq:S}--\eqref{eq:kappamodel}, so this reduction is accounted for. The Gamma family is
positive, has one shape parameter, contains the exponential distribution ($e=3$) and tends to the
constant 1 as $e\to0$. The depth is capped (30 by default) to avoid unrealistic peaks.

\subsection{Burst width and duration from the record}\label{sec:structure}

The program estimates $B$ and $T_{c,j}$ from the reference. The record is cut into frames of 0.5\,s
(hop 0.25\,s, Hann window, i.e.\ the taper $\sin^2(\pi n/L)$ over the $L$ samples of a frame); in every
frame the power in the 1/3-octave band around every node is computed, and its logarithm
$\beta_{j\ell}$ is taken. Even for a stationary signal these values scatter: the power of a band
of $n_l$ lines is a sum of $2n_l$ squared Gaussian variables (the real and imaginary parts of each
line), i.e.\ a \emph{chi-square} variable with $\nu = 2n_l$ degrees of freedom up to a scale factor.
Half of such a variable, $Y$, has the Gamma distribution of shape $k = \nu/2$ and scale 1 (density
$y^{k-1}e^{-y}/\Gamma(k)$), whose moments are $\E[Y^s] = \int_0^\infty y^{s+k-1}e^{-y}\dd y/\Gamma(k) =
\Gamma(k+s)/\Gamma(k)$. The logarithm of $\E[e^{s\ln Y}] = \E[Y^s]$, the cumulant generating function
of $\ln Y$, is therefore $\ln\Gamma(k+s)-\ln\Gamma(k)$; its second derivative at $s=0$ is the
variance of $\ln Y$, the \emph{trigamma function} $\psi'(k) = \frac{\dd^2}{\dd k^2}\ln\Gamma(k)$ at
$k=\nu/2$ (a scale factor does not change the variance of a logarithm). This noise variance is subtracted from the variance of each band before the
correlations are normalised (the effective $\nu$ of a band with $n_l$ lines is taken as $2n_l/1.5$,
since the Hann window correlates neighbouring lines). Bands whose remaining variance is negligible
do not fluctuate and are left out. Then:
\begin{itemize}[nosep]
\item the time correlation of each band is its normalised autocorrelation; the lag where it falls
below $1/e$, halved, is $T_{c,j}$ (smoothed along the nodes and interpolated to all of them);
\item the frequency correlation is the average correlation between bands as a function of their
distance in octaves; the distance where it falls below $1/e$, halved, is $B$ (between 1/12 and 3
octaves).
\end{itemize}
The burst duration is finally kept above $\max(2\tau_i, 0.25\,\mathrm{s})$: shorter bursts would not
reach the responses (Section~\ref{sec:papoulis}) and would only require deeper modulation.

\subsection{Keeping the time--frequency pattern of the record}\label{sec:pattern}

Instead of a random field with the statistics of the record, the field can follow the record itself
(\code{modulation = "pattern"}): the bursts then occur in the same bands and in the same order as in
the reference. The logarithm of the band power of the record, $\beta_j(t)$, is computed at every node
with the frame length that the synthesis uses at that node (so the time resolution is the one the
synthesis can reproduce), smoothed lightly, interpolated on the frame grid of the test and centred,
$\tilde\beta_j = \beta_j - \tavg{\beta_j}$. The field is
\begin{equation}
  m_j(t) = \frac{\exp\bigl(\gamma_j\tilde\beta_j(t)\bigr)}{\tavg{\exp(\gamma_j\tilde\beta_j)}} ,
\end{equation}
where the exponent $\gamma_j\ge0$ sets the depth: $\gamma_j=1$ reproduces the depth of the record, and
the design is free to change it to meet $\kappa_t$ and $D_t$. The map from $\gamma_j$ to the excess
$e_j = 3\Var(m_j)$ is tabulated on a grid of exponents and made monotone, so that the design can
work with $e_j$ as for the random field.

A test shorter than the record needs a \emph{mapping} of the reference time onto the test time:
\begin{itemize}[nosep]
\item \code{blocks} (default): blocks of a few seconds taken evenly across the whole record and
concatenated; every part of the record is represented and the burst durations are kept;
\item \code{compress}: the whole record squeezed into the test (the band power is averaged over each
test frame); the order of all events is kept and the bursts are shortened;
\item a test as long as the record uses it frame by frame (\code{identity}); a longer one repeats it
(\code{loop}).
\end{itemize}
The agreement between the band-power pattern of the result and the mapped reference is reported as
the median, over the nodes, of their correlation coefficient (1: same pattern).
